\subsection{Problematyka i zakres pracy}\label{section_problem}

% \quad We współczesnych środowiskach zawodowych oraz akademickich regularnie obserwuje się wzrost oczekiwań dotyczących liczby i złożoności realizowanych zadań. Pomimo rosnącej presji efektywnościowej, wiele osób nie dysponuje narzędziami ani umiejętnościami pozwalającymi na skuteczne zarządzanie własnym czasem. Prowadzi to do chronicznego przebodźcowania, spadku efektywności, a w skrajnych przypadkach – do wypalenia zawodowego. Jednym z czynników wpływających na nieefektywne planowanie jest nieuwzględnianie zmiennej natury funkcjonowania człowieka, w tym zmian wydajnosci związanych z rytmem dobowym i koniecznością odpoczynku. W efekcie powielane są niekorzystne nawyki, zamiast świadomego wykorzystania wiedzy o własnym rytmie dobowym do optymalizacji pracy.
\quad We współczesnych środowiskach zawodowych oraz akademickich regularnie obserwuje się wzrost oczekiwań dotyczących liczby i złożoności realizowanych zadań. Jednocześnie skuteczne planowanie aktywności człowieka stanowi złożony problem, ponieważ jego wydajność nie jest wartością stałą i może zmieniać się w zależności od czynników wewnętrznych i zewnętrznych. W szczególności znaczenie mają zmęczenie, potrzeba odpoczynku, rytm dobowy oraz indywidualne różnice w sposobie funkcjonowania. Nieuwzględnianie tej zmienności może prowadzić do tworzenia planów niedostosowanych do rzeczywistych możliwości użytkownika, a w konsekwencji do spadku efektywności, narastania obciążenia oraz powielania negatywnych nawyków.
% \\\indent Odpowiedzią na te wyzwania mogą być inteligentne systemy wspomagania decyzji, automatyzujące proces harmonogramowania w sposób spersonalizowany. Nowoczesne rozwiązania tego typu nie ograniczają się jedynie do sztywnych parametrów samych zadań (takich jak terminy czy priorytety). Opierają one swoje decyzje na modelowaniu zachowań użytkownika, analizie informacji zwrotnej o faktycznej użyteczności planu oraz iteracyjnej adaptacji algorytmów w czasie. Pozwala to na dostosowywanie systemu do indywidualnych cech użytkownika oraz obserwowanej zmienności jego wydajności w czasie.
\\\indent Odpowiedzią na te wyzwania mogą być inteligentne systemy wspomagania decyzji, automatyzujące proces harmonogramowania w sposób spersonalizowany. W przeciwieństwie do klasycznych metod opartych na sztywnych parametrach i wcześniej zdefiniowanych regułach, rozwiązania adaptacyjne mogą uwzględniać zmiany stanu użytkownika, jego historyczną wydajność oraz informacje zwrotne dotyczące realizacji planu. Pozwala to na bieżące dostosowywanie harmonogramu do indywidualnych cech użytkownika oraz zmienności jego funkcjonowania w czasie.
% \\\indent Choć metody automatycznego planowania są szeroko rozwijane w kontekście procesów systemowych czy interakcji człowiek-maszyna w przemyśle, istnieje wyraźna luka w ich adaptacji do dynamicznego i zróżnicowanego środowiska, jakim jest codzienne życie człowieka. 
\\\indent Choć metody automatycznego planowania są szeroko rozwijane w kontekście procesów systemowych, produkcyjnych czy współpracy człowieka z maszyną, ich zastosowanie w codziennym, indywidualnym planowaniu zadań nadal pozostaje ograniczone. Istotnym problemem jest również sposób reprezentowania czynników ludzkich, takich jak zmęczenie i regeneracja. W części istniejących rozwiązań odpoczynek traktowany jest jako aktywny element planu, jednak w wielu przypadkach jego wpływ na stan użytkownika nie jest jawnie modelowany.
% \\\indent Niniejsza praca podejmuje tę tematykę poprzez formalizację wpływu zmęczenia, rytmu dobowego oraz przerw na procesy decyzyjne. Zakres pracy obejmuje przegląd literatury i identyfikację metod optymalizacyjnych oraz podejść opartych na sztucznej inteligencji. W ramach badań opracowany i zaimplementowany zostanie jednolity model problemu planowania, który pozwoli na przeprowadzenie eksperymentów porównawczych wybranych algorytmów pod kątem efektywności oraz zgodności harmonogramu z przyjętymi ograniczeniami psychofizycznymi użytkownika.
\\\indent Niniejsza praca podejmuje tę tematykę poprzez formalizację wpływu rytmu dobowego, chronotypu, narastającego zmęczenia oraz regeneracji podczas odpoczynku na proces planowania zadań. Zakres pracy obejmuje analizę istniejących metod automatycznego harmonogramowania oraz opracowanie modelu symulacyjnego użytkownika, umożliwiającego prowadzenie eksperymentów porównawczych dla wybranych algorytmów w jednolitych warunkach.

% \newpage
\subsection{Cele pracy}\label{section_goals}

% \quad Głównym celem pracy jest opracowanie oraz eksperymentalna ocena wybranych metod automatycznego planowania zadań, z ukierunkowaniem na personalizację harmonogramu dla konkretnego użytkownika. W szczególności analizowany jest wpływ rytmu dobowego, zmienność funkcjonowania użytkownika w czasie pod wpływem narastającego zmęczenia i regeneracji w czasie odpoczynku. W związku z tym kluczowe jest zapoznanie się z metodami modelowania poziomów energii oraz wydajności człowieka z uwzględnieniem różnych chronotypów, a także dobór inteligentnych metod harmonogramowania umożliwiających operowanie na wielu zmiennych.
% \\\indent W celu zapewnienia jednolitych warunków porównania opracowany zostanie model symulacyjny użytkownika umożliwiający reprezentację różnych profili funkcjonowania oraz ich reakcji na wykonywanie zadań i odpoczynek. Na tej podstawie przeprowadzone zostanie porównanie modeli planistycznych dedykowanych dla jednostek z wykorzystaniem zdefiniowanych scenariuszy i miar jakośći. Eksperyment zostanie zrealizowany na bazie augmentowanych, ogólnodostępnych danych, co zapewni uproszczenie oraz weryfikowalność całego procesu.
% \\\indent Dodatkowym celem jest ocena wyników uzyskanych przez wybrane metody planowania pod kątem ich dostosowania do dynamicznego środowiska, potrzeb konkretnych użytkowników, narzuconych ograniczeń zadań oraz reakcji na zakłócenia, w oparciu o zdefiniowane miary efektywności.
\quad Głównym celem pracy jest opracowanie oraz eksperymentalna ocena wybranych metod automatycznego planowania zadań, z ukierunkowaniem na personalizację harmonogramu dla konkretnego użytkownika. W szczególności analizowany jest wpływ rytmu dobowego oraz zmienności funkcjonowania użytkownika w czasie, wynikającej z narastającego zmęczenia i regeneracji zachodzącej podczas odpoczynku. W związku z tym kluczowe jest zapoznanie się z metodami modelowania poziomu energii i wydajności człowieka z uwzględnieniem różnych chronotypów, a także dobór inteligentnych metod harmonogramowania umożliwiających uwzględnienie wielu zmiennych opisujących zarówno użytkownika, jak i realizowane zadania.
\\\indent W celu zapewnienia jednolitych warunków porównania opracowany zostanie model symulacyjny użytkownika, umożliwiający reprezentację różnych profili funkcjonowania oraz ich reakcji na wykonywanie zadań i odpoczynek. Na tej podstawie przeprowadzone zostanie porównanie wybranych modeli planistycznych dostosowanych do poszczególnych profili użytkowników, z wykorzystaniem zdefiniowanych scenariuszy badawczych oraz miar jakości. Eksperyment zostanie zrealizowany na bazie augmentowanych, ogólnodostępnych danych, co pozwoli na uproszczenie przygotowania środowiska badawczego oraz zapewnienie weryfikowalności całego procesu.
\\\indent Dodatkowym celem pracy jest ocena wyników uzyskanych przez wybrane metody planowania pod kątem ich zdolności do dostosowania się do dynamicznego środowiska, potrzeb konkretnych użytkowników, narzuconych ograniczeń związanych z zadaniami oraz pojawiających się zakłóceń. Ocena zostanie przeprowadzona w oparciu o zdefiniowane miary efektywności, pozwalające na porównanie jakości generowanych harmonogramów oraz stopnia ich dopasowania do przyjętego modelu funkcjonowania użytkownika.